\chapter{Week 8: Holomorphs and Primitive Actions}

\section{Regular Actions}

\begin{lemma}
The subgroups \(\widehat X\) and \(\check X\) are regular subgroups of
\(\operatorname{Sym}(X)\), and they centralize each other.
\end{lemma}

\begin{proof}
The action of \(\widehat X\) is regular because for any \(x,y\in X\), the unique
element sending \(x\) to \(y\) is
\[
    \widehat{x^{-1}y}.
\]
Similarly, the action of \(\check X\) is regular because the unique element sending
\(x\) to \(y\) is
\[
    \check{yx^{-1}}.
\]

Finally, for all \(a,b,x\in X\),
\[
    x^{\widehat a\check b}
    =
    b^{-1}xa
    =
    x^{\check b\widehat a}.
\]
Hence every element of \(\widehat X\) commutes with every element of
\(\check X\).
\end{proof}

\begin{definition}[Equivalent permutation actions]
Let \(G\) act on two sets \(\Omega\) and \(\Delta\). These two actions are called
\textbf{equivalent} if there exists a bijection
\[
    \varphi:\Omega\to\Delta
\]
such that
\[
    \varphi(\omega^g)=\varphi(\omega)^g
    \qquad
    \text{for all } \omega\in\Omega,\ g\in G.
\]
\end{definition}

\begin{corollary}
Every regular action of a group \(X\) is equivalent to the action of
\(\widehat X\) on \(X\), and also to the action of \(\check X\) on \(X\).
\end{corollary}

\begin{proof}
Let \(X\) act regularly on \(\Omega\). Fix \(\omega\in\Omega\).

First define
\[
    \varphi:X\to\Omega,
    \qquad
    a\mapsto \omega^a.
\]
Since the action is regular, \(\varphi\) is a bijection. For \(a,b\in X\), we have
\[
    \varphi(a^{\widehat b})
    =
    \varphi(ab)
    =
    \omega^{ab}
    =
    (\omega^a)^b
    =
    \varphi(a)^b.
\]
Thus the given regular action is equivalent to the right regular action.

For the left regular action, define
\[
    \psi:X\to\Omega,
    \qquad
    a\mapsto \omega^{a^{-1}}.
\]
Again \(\psi\) is a bijection. For \(a,b\in X\),
\[
    \psi(a^{\check b})
    =
    \psi(b^{-1}a)
    =
    \omega^{(b^{-1}a)^{-1}}
    =
    \omega^{a^{-1}b}
    =
    (\omega^{a^{-1}})^b
    =
    \psi(a)^b.
\]
Hence the given regular action is also equivalent to the left regular action.
\end{proof}

\begin{lemma}[Centralizer of a regular group]
Let \(X\) act on itself by right multiplication. Then
\[
    C_{\operatorname{Sym}(X)}(\widehat X)=\check X.
\]
\end{lemma}

\begin{proof}
We already know that \(\check X\leq C_{\operatorname{Sym}(X)}(\widehat X)\).

Conversely, let \(c\in C_{\operatorname{Sym}(X)}(\widehat X)\). Put
\[
    1^c=a.
\]
Then for every \(x\in X\),
\[
    x^c
    =
    (1^{\widehat x})^c
    =
    (1^c)^{\widehat x}
    =
    ax.
\]
But the map \(x\mapsto ax\) is \(\check{a^{-1}}\). Hence \(c\in\check X\).
Therefore
\[
    C_{\operatorname{Sym}(X)}(\widehat X)=\check X.
\]
\end{proof}

\section{The Holomorph}

\begin{definition}[Holomorph]
Let \(G\) be a group. The \textbf{holomorph} of \(G\) is
\[
    \operatorname{Hol}(G)
    =
    G:\operatorname{Aut}(G).
\]
Equivalently, under the right action convention, we may realize it as
\[
    \operatorname{Hol}(G)
    =
    \widehat G\operatorname{Aut}(G)
    \leq
    \operatorname{Sym}(G).
\]

As an abstract semidirect product,
\[
    \operatorname{Hol}(G)
    =
    \{(g,\alpha)\mid g\in G,\ \alpha\in\operatorname{Aut}(G)\},
\]
with multiplication
\[
    (g,\alpha)(h,\beta)
    =
    \bigl(g h^{\alpha^{-1}},\alpha\beta\bigr).
\]
\end{definition}

\begin{lemma}[Normalizer of a regular subgroup]
Let \(M\leq \operatorname{Sym}(\Omega)\) be regular. After identifying
\(\Omega\) with \(M\), we have
\[
    N_{\operatorname{Sym}(\Omega)}(M)
    =
    M:\operatorname{Aut}(M)
    =
    \operatorname{Hol}(M).
\]
\end{lemma}

\begin{proof}
Identify \(\Omega\) with \(M\), so that \(M=\widehat M\) acts by right
multiplication.

Let \(G=N_{\operatorname{Sym}(M)}(\widehat M)\), and let \(1\in M\) be the
identity element. Since \(\widehat M\) is regular, we have
\[
    G=\widehat M G_1,
\]
where \(G_1\) is the stabilizer of \(1\).

We claim that
\[
    G_1\leq \operatorname{Aut}(M).
\]
Let \(b\in G_1\). Since \(b\) normalizes \(\widehat M\), for every \(m\in M\)
there exists \(m'\in M\) such that
\[
    b^{-1}\widehat m b=\widehat{m'}.
\]
Evaluating both sides at \(1\), we get
\[
    m'=m^b.
\]
Hence
\[
    b^{-1}\widehat m b=\widehat{m^b}.
\]
Therefore, for \(m,n\in M\),
\[
    (mn)^b
    =
    \left(1^{\widehat m\widehat n}\right)^b
    =
    1^{b^{-1}\widehat m\widehat n b}
    =
    1^{(b^{-1}\widehat m b)(b^{-1}\widehat n b)}
    =
    1^{\widehat{m^b}\widehat{n^b}}
    =
    m^b n^b.
\]
Thus \(b\) is an automorphism of \(M\).

Conversely, every automorphism of \(M\) normalizes the right regular subgroup
\(\widehat M\). Hence
\[
    N_{\operatorname{Sym}(M)}(\widehat M)
    =
    \widehat M:\operatorname{Aut}(M).
\]
\end{proof}

\section{Holomorph Simple Type}

\begin{theorem}[Holomorph simple type]
Let \(G\leq \operatorname{Sym}(\Omega)\). Suppose
\[
    M\times N\unlhd G
\]
and suppose that \(M\) and \(N\) are regular on \(\Omega\), with
\[
    M\cong N\cong T,
\]
where \(T\) is a nonabelian simple group.

After identifying \(\Omega\) with \(M\), we have
\[
    M:\operatorname{Inn}(M)
    \leq
    G
    \leq
    M:\operatorname{Aut}(M).
\]
\end{theorem}

\begin{proof}
Since \(M\) is regular, we identify \(\Omega\) with \(M\), so that
\[
    M=\widehat M.
\]
Since \(M\unlhd G\), the group \(G\) normalizes \(M\). Hence, by the normalizer
lemma,
\[
    G\leq N_{\operatorname{Sym}(M)}(\widehat M)
    =
    M:\operatorname{Aut}(M).
\]

It remains to identify the lower bound. Since \(M\times N\unlhd G\), the
subgroups \(M\) and \(N\) commute. Thus
\[
    N\leq C_{\operatorname{Sym}(M)}(\widehat M).
\]
By the centralizer lemma,
\[
    C_{\operatorname{Sym}(M)}(\widehat M)=\check M.
\]
Since \(N\) and \(\check M\) are both regular and have the same order, we get
\[
    N=\check M.
\]

Now \(M\cong T\) is centerless. Hence
\[
    \widehat M\cap \check M=1.
\]
The point stabilizer of \(1\in M\) in \(\widehat M\check M\) consists of the
elements
\[
    \widehat m\check m
    \qquad (m\in M),
\]
because
\[
    1^{\widehat m\check m}=m^{-1}m=1.
\]
Moreover,
\[
    x^{\widehat m\check m}
    =
    m^{-1}xm,
\]
so this point stabilizer is naturally identified with
\[
    \operatorname{Inn}(M).
\]
Therefore
\[
    \widehat M\check M
    =
    M:\operatorname{Inn}(M)
    \leq G.
\]
Thus
\[
    M:\operatorname{Inn}(M)
    \leq
    G
    \leq
    M:\operatorname{Aut}(M).
\]
\end{proof}

\section{Holomorph Compound Type}

\begin{theorem}[Holomorph compound type]
Let \(G\leq \operatorname{Sym}(\Omega)\). Suppose
\[
    M\times N\unlhd G,
\]
where \(M\) and \(N\) are regular on \(\Omega\), and
\[
    M\cong N\cong T^\ell
\]
with \(\ell\geq 2\) and \(T\) a nonabelian simple group.

After identifying \(\Omega\) with \(M\), we have
\[
    M:\operatorname{Inn}(M)
    \leq
    G
    \leq
    M:\operatorname{Aut}(M).
\]
Moreover,
\[
    \operatorname{Aut}(M)
    \cong
    \operatorname{Aut}(T)\wr S_\ell.
\]
Thus
\[
    M:\operatorname{Inn}(M)
    \leq
    G
    \leq
    M:\bigl(\operatorname{Aut}(T)\wr S_\ell\bigr).
\]
\end{theorem}

\begin{proof}
The proof is the same as in the holomorph simple case.

Identify \(\Omega\) with \(M\), so \(M=\widehat M\). Since \(M\unlhd G\),
\[
    G\leq N_{\operatorname{Sym}(M)}(\widehat M)
    =
    M:\operatorname{Aut}(M).
\]
Since \(M\) and \(N\) commute and \(N\) is regular,
\[
    N\leq C_{\operatorname{Sym}(M)}(\widehat M)=\check M.
\]
Hence \(N=\check M\). Since \(T^\ell\) is centerless, we have
\[
    \widehat M\cap \check M=1.
\]
The point stabilizer of \(1\in M\) in \(\widehat M\check M\) is identified with
\(\operatorname{Inn}(M)\). Therefore
\[
    \widehat M\check M
    =
    M:\operatorname{Inn}(M)
    \leq G.
\]

Finally, since
\[
    M=T_1\times \cdots \times T_\ell
    \qquad
    (T_i\cong T),
\]
the automorphism group of \(M\) permutes the simple direct factors and acts on
each factor by an automorphism of \(T\). Hence
\[
    \operatorname{Aut}(M)
    \cong
    \operatorname{Aut}(T)\wr S_\ell.
\]
\end{proof}

\begin{remark}
In the quasiprimitive holomorph compound case, the group \(G\) induces a
transitive action on the set of simple direct factors
\[
    \{T_1,\dots,T_\ell\}.
\]
Indeed, if the induced action had more than one orbit, then the product of the
factors in one orbit would give a nontrivial proper normal subgroup of \(G\)
inside \(M\), contradicting the minimal normality of \(M\).
\end{remark}

\section{The Frattini Argument for Permutation Groups}

\begin{lemma}[Frattini argument for permutation groups]
Let \(H\leq G\leq \operatorname{Sym}(\Omega)\), and assume that \(G\) is
transitive on \(\Omega\). Fix \(\omega\in\Omega\). Then \(H\) is transitive on
\(\Omega\) if and only if
\[
    G=G_\omega H.
\]
In particular, if \(H\unlhd G\), then this is equivalent to
\[
    G=HG_\omega.
\]
\end{lemma}

\begin{proof}
Suppose first that \(H\) is transitive. Let \(g\in G\). Since \(H\) is transitive,
there exists \(h\in H\) such that
\[
    \omega^h=\omega^g.
\]
Then
\[
    \omega^{gh^{-1}}
    =
    (\omega^g)^{h^{-1}}
    =
    (\omega^h)^{h^{-1}}
    =
    \omega.
\]
Hence \(gh^{-1}\in G_\omega\), and therefore
\[
    g\in G_\omega H.
\]
Thus \(G=G_\omega H\).

Conversely, suppose \(G=G_\omega H\). For every \(g\in G\), write
\[
    g=kh
    \qquad
    (k\in G_\omega,\ h\in H).
\]
Then
\[
    \omega^g
    =
    \omega^{kh}
    =
    (\omega^k)^h
    =
    \omega^h.
\]
Thus every point in the \(G\)-orbit of \(\omega\) lies in the \(H\)-orbit of
\(\omega\). Since \(G\) is transitive, \(H\) is transitive.
\end{proof}

\begin{lemma}
Let \(M\unlhd G\leq \operatorname{Sym}(\Omega)\), and suppose that \(M\) is
regular on \(\Omega\). Then
\[
    G=M:G_\omega
\]
for every \(\omega\in\Omega\).
\end{lemma}

\begin{proof}
Since \(M\) is transitive on \(\Omega\), the Frattini argument gives
\[
    G=G_\omega M.
\]
Since \(M\unlhd G\), we also have
\[
    G=MG_\omega.
\]
Moreover,
\[
    M\cap G_\omega=1
\]
because \(M\) is regular. Hence
\[
    G=M:G_\omega.
\]
\end{proof}

\section{Blocks and Primitive Actions}

\begin{definition}[Block]
Let \(G\) be a transitive permutation group on \(\Omega\). A subset
\(B\subseteq \Omega\) is called a \textbf{block} for \(G\) if for every
\(g\in G\),
\[
    B^g=B
    \qquad
    \text{or}
    \qquad
    B^g\cap B=\varnothing.
\]
\end{definition}

\begin{definition}[Block system]
Let \(B\) be a block for \(G\). The set
\[
    \mathcal B=\{B^g\mid g\in G\}
\]
is called the \textbf{block system} generated by \(B\).
\end{definition}

\begin{proposition}[Basic properties of block systems]
Let \(G\) be transitive on \(\Omega\), and let
\[
    \mathcal B=\{B_1,B_2,\dots,B_m\}
\]
be a block system. Then:
\begin{enumerate}[label=(\arabic*)]
    \item each \(B_i\) is a block;
    \item \(B_i\cap B_j=\varnothing\) whenever \(i\neq j\);
    \item \(|B_1|=\cdots=|B_m|\);
    \item \(G\) induces a transitive action on \(\mathcal B\).
\end{enumerate}
\end{proposition}

\begin{proof}
Each \(B_i\) has the form \(B^g\) for some \(g\in G\), and conjugating the block
condition shows that \(B^g\) is again a block. Hence each \(B_i\) is a block.

If \(B_i\cap B_j\neq\varnothing\), then by the block condition they must be equal.
Thus distinct blocks in \(\mathcal B\) are disjoint.

All blocks in \(\mathcal B\) are images of \(B\) under permutations, so they have the
same cardinality.

Finally, the action of \(G\) on \(\mathcal B\) is transitive by construction, since every
block in \(\mathcal B\) is of the form \(B^g\).
\end{proof}

\begin{definition}[Primitive and imprimitive actions]
Let \(G\) be transitive on \(\Omega\).

If there exists a block \(B\subseteq \Omega\) such that
\[
    1<|B|<|\Omega|,
\]
then \(G\) is called \textbf{imprimitive} on \(\Omega\).

If no such block exists, then \(G\) is called \textbf{primitive} on \(\Omega\).
\end{definition}

\begin{lemma}[Primitive action and maximal point stabilizer]
Let \(G\leq \operatorname{Sym}(\Omega)\) be transitive, and fix
\(\omega\in\Omega\). Then \(G\) is primitive on \(\Omega\) if and only if
\(G_\omega\) is a maximal subgroup of \(G\).
\end{lemma}

\begin{proof}
Suppose first that \(G\) is imprimitive. Let \(B\) be a nontrivial block with
\[
    \omega\in B.
\]
Let
\[
    G_B=\{g\in G\mid B^g=B\}
\]
be the setwise stabilizer of \(B\). Then
\[
    G_\omega\leq G_B.
\]
Since \(B\) is nontrivial, choose \(\omega'\in B\) with \(\omega'\neq \omega\).
By transitivity, there exists \(g\in G\) such that
\[
    \omega^g=\omega'.
\]
Then \(B^g\cap B\neq\varnothing\), so \(B^g=B\). Hence
\[
    g\in G_B.
\]
But \(g\notin G_\omega\), so
\[
    G_\omega<G_B.
\]
Also \(G_B<G\), since otherwise \(B\) would be \(G\)-invariant, and transitivity
would force \(B=\Omega\). Therefore \(G_\omega\) is not maximal.

Conversely, suppose that \(G_\omega\) is not maximal. Then there exists a subgroup
\(H\) such that
\[
    G_\omega<H<G.
\]
Define
\[
    B=\omega^H.
\]
We claim that \(B\) is a block. Let \(g\in G\), and suppose that
\[
    B^g\cap B\neq\varnothing.
\]
Then there exist \(h_1,h_2\in H\) such that
\[
    \omega^{h_1g}=\omega^{h_2}.
\]
Thus
\[
    \omega^{h_1gh_2^{-1}}=\omega,
\]
so
\[
    h_1gh_2^{-1}\in G_\omega\leq H.
\]
Hence \(g\in H\), and therefore
\[
    B^g=B.
\]
Thus \(B\) is a block.

Since \(H>G_\omega\), the block \(B\) has more than one point. Since \(H<G\),
the subgroup \(H\) is not transitive; otherwise the Frattini argument would imply
\(G=G_\omega H=H\), a contradiction. Hence \(B\neq \Omega\). Therefore \(B\)
is a nontrivial block, and \(G\) is imprimitive.
\end{proof}

\section{Quasiprimitive Groups}

\begin{definition}[Quasiprimitive permutation group]
Let \(G\leq \operatorname{Sym}(\Omega)\) be transitive. We say that \(G\) is
\textbf{quasiprimitive} on \(\Omega\) if every nontrivial normal subgroup of \(G\)
is transitive on \(\Omega\).
\end{definition}

\begin{lemma}[Blocks for groups with a regular normal subgroup]
Let \(G=M:G_\omega\leq \operatorname{Sym}(\Omega)\), where \(M\unlhd G\) is
regular on \(\Omega\). Identify \(\Omega\) with \(M\), and identify \(\omega\) with
the identity element \(1\in M\).

Then the blocks of \(G\) containing \(1\) are exactly the subgroups
\(K\leq M\) that are invariant under \(G_\omega\).
\end{lemma}

\begin{proof}
Let \(B\) be a block containing \(1\). Since \(M\) is regular, we may regard
\(B\) as a subset of \(M\).

We first show that \(B\) is a subgroup of \(M\). Let \(a,b\in B\). Since
\(b\in B\), the translate
\[
    B^{\widehat{b^{-1}}}
\]
contains \(1\). Hence
\[
    B^{\widehat{b^{-1}}}\cap B\neq\varnothing.
\]
Because \(B\) is a block, we have
\[
    B^{\widehat{b^{-1}}}=B.
\]
Therefore
\[
    ab^{-1}\in B.
\]
By the subgroup criterion, \(B\leq M\).

Next let \(x\in G_\omega\). Since \(x\) fixes \(1\), the block \(B^x\) also contains
\(1\). Hence
\[
    B^x\cap B\neq\varnothing.
\]
Thus
\[
    B^x=B.
\]
So \(B\) is \(G_\omega\)-invariant.

Conversely, suppose \(K\leq M\) is \(G_\omega\)-invariant. Let \(g\in G\). Since
\[
    G=M:G_\omega,
\]
we may write
\[
    g=\widehat m x
    \qquad
    (m\in M,\ x\in G_\omega).
\]
Then
\[
    K^g
    =
    K^{\widehat m x}
    =
    (Km)^x
    =
    K^x m^x
    =
    K m^x,
\]
because \(K^x=K\). Hence \(K^g\) is a right coset of \(K\) in \(M\). Two right
cosets of \(K\) are either equal or disjoint, so
\[
    K^g=K
    \qquad
    \text{or}
    \qquad
    K^g\cap K=\varnothing.
\]
Thus \(K\) is a block.
\end{proof}

\begin{theorem}
A quasiprimitive permutation group of type \(\mathrm{HA}\), \(\mathrm{HS}\), or
\(\mathrm{HC}\) is primitive.
\end{theorem}

\begin{proof}
In each of the types \(\mathrm{HA}\), \(\mathrm{HS}\), and \(\mathrm{HC}\), the group
\(G\) has a regular minimal normal subgroup \(M\). Hence
\[
    G=M:G_\omega.
\]
By the previous lemma, if \(G\) were imprimitive, then there would exist a
nontrivial proper subgroup
\[
    1<K<M
\]
that is invariant under \(G_\omega\).

We show that this is impossible in each case.

\medskip

\noindent
\textbf{Case \(\mathrm{HA}\).}
Here \(M\) is abelian and is the unique minimal normal subgroup of \(G\). Since
\(K\leq M\) and \(M\) is abelian, \(M\) normalizes \(K\). Since \(K\) is
\(G_\omega\)-invariant, \(G_\omega\) also normalizes \(K\). Thus
\[
    K\unlhd M:G_\omega=G.
\]
This contradicts the minimal normality of \(M\), because
\[
    1<K<M.
\]
Therefore no such \(K\) exists, and \(G\) is primitive.

\medskip

\noindent
\textbf{Case \(\mathrm{HS}\).}
Here
\[
    M\cong T
\]
for a nonabelian simple group \(T\), and from the holomorph simple structure,
\[
    \operatorname{Inn}(M)\leq G_\omega.
\]
If \(K\leq M\) is \(G_\omega\)-invariant, then in particular \(K\) is invariant
under \(\operatorname{Inn}(M)\). Hence
\[
    K\unlhd M.
\]
Since \(M\) is simple, we get
\[
    K=1
    \qquad
    \text{or}
    \qquad
    K=M.
\]
Thus no nontrivial proper \(K\) exists, and \(G\) is primitive.

\medskip

\noindent
\textbf{Case \(\mathrm{HC}\).}
Here
\[
    M=T_1\times \cdots \times T_\ell,
    \qquad
    T_i\cong T,
\]
where \(T\) is nonabelian simple and \(\ell\geq 2\). Again
\[
    \operatorname{Inn}(M)\leq G_\omega.
\]
Thus a \(G_\omega\)-invariant subgroup \(K\leq M\) is normalized by
\(\operatorname{Inn}(M)\), hence is normal in \(M\). Therefore \(K\) is a product
of some of the simple direct factors \(T_i\).

But in the holomorph compound case, \(G_\omega\) acts transitively on the set
\[
    \{T_1,\dots,T_\ell\}.
\]
Therefore the set of direct factors appearing in \(K\) is either empty or all of
\(\{T_1,\dots,T_\ell\}\). Hence
\[
    K=1
    \qquad
    \text{or}
    \qquad
    K=M.
\]
So no nontrivial proper \(G_\omega\)-invariant subgroup \(K\) exists. Therefore
\(G\) is primitive.
\end{proof}

\begin{proposition}
Every primitive permutation group is quasiprimitive.
\end{proposition}

\begin{proof}
Let \(G\leq \operatorname{Sym}(\Omega)\) be primitive, and let
\[
    1\neq H\unlhd G.
\]
We prove that \(H\) is transitive.

Fix \(\omega\in\Omega\). If
\[
    H\leq G_\omega,
\]
then for every \(g\in G\),
\[
    H^g=H\leq G_{\omega^g}.
\]
Since \(G\) is transitive, the points \(\omega^g\) run over all of \(\Omega\). Hence
\(H\) fixes every point of \(\Omega\). Since \(G\leq \operatorname{Sym}(\Omega)\)
is faithful, this forces
\[
    H=1,
\]
contradicting \(H\neq 1\).

Therefore
\[
    H\nleq G_\omega.
\]
Thus
\[
    G_\omega<HG_\omega\leq G.
\]
Since \(G\) is primitive, \(G_\omega\) is maximal in \(G\). Hence
\[
    HG_\omega=G.
\]
By the Frattini argument, \(H\) is transitive on \(\Omega\). Therefore every
nontrivial normal subgroup of \(G\) is transitive, so \(G\) is quasiprimitive.
\end{proof}

\begin{example}[Quasiprimitive but not primitive]
Let \(T\) be a nonabelian simple group. Let \(\widehat T\) act regularly on the set
\(T\).

Then \(\widehat T\) is quasiprimitive on \(T\). Indeed, since \(T\) is simple, the
only nontrivial normal subgroup of \(\widehat T\) is \(\widehat T\) itself, and
\(\widehat T\) is transitive.

However, \(\widehat T\) is not primitive. The stabilizer of the identity element
\(1\in T\) is
\[
    \widehat T_1=1.
\]
Since \(T\) is nonabelian simple, \(T\) has a nontrivial proper subgroup. Hence
\(1\) is not a maximal subgroup of \(\widehat T\). By the maximal stabilizer
criterion, the regular action of \(\widehat T\) on \(T\) is not primitive.
\end{example}